\section{Learning Markov Networks}
\label{sec:learning-mn}

Learning consists in choosing the parameter vector $w$ of the Markov
Network from a training set
$\mathcal{D} = \{(x^i, y^i)\}_{i=1}^{m} \subset \mathcal{X} \times
\mathcal{Y}^{N}$. The approach adopted in this thesis is the
\emph{max-margin} one due to Taskar et al.\ and Tsochantaridis et
al.~\cite{Taskar2004, Tsochantaridis2005}: $w$ is chosen so that, for
every training pair, the score of the ground-truth labeling exceeds
the score of every wrong labeling by a margin proportional to how
wrong it is. Because exact max-margin training is intractable on
cyclic graphs, the training loss is further replaced by a tractable
LP-relaxed surrogate following Franc \& Yermakov~\cite{franc2021}.

\subsection{Max-margin formulation}
\label{subsec:max-margin}

Assuming the unary and pairwise potentials of
Section~\ref{sec:markov-networks} are linear in the parameters
$w \in \mathbb{R}^{d}$, the MAP predictor of
Eq.~\eqref{eq:map-prediction} reduces to the linear classifier
\begin{equation}
    h(x; w) \;=\; \argmax_{y \in \mathcal{Y}^{N}}\,
    \langle w, \psi(x, y) \rangle,
    \label{eq:linear-classifier}
\end{equation}
with the joint feature map $\psi$ defined in
Eq.~\eqref{eq:joint-feature}. Given a task-specific loss
$\ell \colon \mathcal{Y}^{N} \times \mathcal{Y}^{N} \to \mathbb{R}_{+}$
(the normalized Hamming loss of
Section~\ref{subsec:loss-functions} in this thesis) and a
regularization constant $\lambda \geq 0$, the
\emph{max-margin Markov network} (M\textsuperscript{3}N) framework
seeks $w$ such that for every training pair the score
$\langle w, \psi(x, y)\rangle$ of the ground truth exceeds the
score of every wrong labeling $y'$ by at least $\ell(y, y')$.
Violations of this margin condition are penalized by the
margin-rescaling loss $\Delta$, leading to the convex unconstrained
problem
\begin{equation}
    w^{*} \;=\; \argmin_{w \in \mathbb{R}^d}\!\left[
    \frac{\lambda}{2}\, \|w\|^{2}
    \;+\; \frac{1}{m}\, \sum_{i=1}^{m}
    \Delta(x^i, y^i, w)
    \right],
    \label{eq:m3n-objective}
\end{equation}
where
\begin{equation}
    \Delta(x, y, w) \;=\;
    \max_{y' \in \mathcal{Y}^{N}}\!
    \Bigl[
        \ell(y, y') + \langle w, \psi(x, y') \rangle
    \Bigr]
    \;-\; \langle w, \psi(x, y) \rangle.
    \label{eq:margin-rescaling}
\end{equation}

$\Delta$ is convex in $w$, as a maximum of affine functions minus
an affine term, and upper-bounds the task loss
$\ell\bigl(y, h(x; w)\bigr)$ achieved by the predictor of
Eq.~\eqref{eq:linear-classifier}; the inner product
$\langle w, \psi(x, y)\rangle$ is the joint score $f(x, y)$ of
Eq.~\eqref{eq:linear-score}. Objective~\eqref{eq:m3n-objective} is
therefore a convex relaxation of regularized empirical risk.

The constant $\lambda$ trades off data fit against the regularizer
$\|w\|^{2}$: a larger $\lambda$ shrinks $w$ toward the origin and
biases the predictor toward simpler decision boundaries, reducing
the risk of overfitting; a smaller $\lambda$ allows $w$ to fit the
training data more aggressively. Its value is selected on a
validation split.

\subsection{Loss-augmented inference}

Both evaluating the loss $\Delta(x, y, w)$ and computing a
subgradient of $\Delta$ with respect to $w$ require solving the
inner maximum in Eq.~\eqref{eq:margin-rescaling}, called the
\emph{loss-augmented inference} problem
\begin{equation}
    y^{\mathrm{LA}}
    \;=\; \argmax_{y' \in \mathcal{Y}^{N}}\!
    \Bigl[\, \ell(y, y') + \langle w, \psi(x, y') \rangle \,\Bigr].
    \label{eq:loss-augmented}
\end{equation}
The normalized Hamming loss is \emph{decomposable}, i.e.\ it
factorizes over output components as
$\ell(y, y') = \frac{1}{N}\sum_{v} \mathbb{I}[y_{v} \neq y'_{v}]$.
The loss term therefore adds a position-wise bonus of $1/N$ to
every label that disagrees with the ground truth, so
Eq.~\eqref{eq:loss-augmented} reduces to MAP inference with
suitably shifted unary potentials --- explicitly, $f_v(x, a)$ is
replaced by $f_v(x, a) + \frac{1}{N}\mathbb{I}[a \neq y_v]$. The
$1/N$ normalization keeps the total task loss in $[0, 1]$
irrespective of $N$ and prevents the augmentation from dominating
the score; without it, a sufficiently confident unary backbone can
close the margin by inflating its outputs by $O(N)$ rather than by
learning useful pairwise structure.

On the linear chain, problem~\eqref{eq:loss-augmented} is an
$O(T K^2)$ Viterbi call, and exact max-margin training of
linear-chain MNs is therefore practical. On the Sudoku graph, in
contrast, loss-augmented inference is NP-hard and a subgradient of
$\Delta$ cannot be computed in polynomial time. This motivates the
LP-relaxed surrogate developed next.

\subsection{Exact training on chains via the structured hinge}
\label{subsec:structured-hinge}

On the linear chain, where loss-augmented inference is exactly
solvable by Viterbi (Section~\ref{subsec:viterbi}), the
M\textsuperscript{3}N objective~\eqref{eq:m3n-objective} can be
optimized directly by stochastic subgradient descent on $\Delta$.
Given a training pair $(x, y)$ and the current parameters $w$,
the oracle returns the most-violating labeling $y^{\mathrm{LA}}$
of Eq.~\eqref{eq:loss-augmented}, and a subgradient of $\Delta$
at $w$ is
\begin{equation}
    \partial_w \Delta(x, y, w)
    \;\ni\; \psi(x, y^{\mathrm{LA}}) - \psi(x, y).
    \label{eq:hinge-subgradient}
\end{equation}
Adding the regularizer contribution $\lambda w$ gives the
subgradient of the full objective~\eqref{eq:m3n-objective}, which
is used to update the parameters via stochastic optimization.

Unlike the LP-relaxed loss developed next, this training path
introduces no auxiliary variables and carries no per-example
state across epochs. In this thesis it is used on the chain
benchmark, where it serves as the reference against which the
LP-relaxed loss is validated for tightness in
Chapter~\ref{chap:experiments}.

\subsection{LP-relaxed margin loss (LP-M3N)}
\label{subsec:lp-m3n}

Following Franc \& Yermakov~\cite{franc2021}, replacing the max in
Eq.~\eqref{eq:loss-augmented} by its LP relaxation
(Section~\ref{subsec:lp-relaxation}) yields an upper bound
$\Delta^{\mathrm{LP}} \geq \Delta$. By LP duality, the relaxed max
admits a reparameterization in which the single global maximum is
replaced by a sum of local maxima mediated by auxiliary variables.
These variables are the duals of the marginal-consistency
constraints of the LP and live on the edges of the graph: for each
edge $\{v, v'\} \in E$ we introduce two vectors $\phi_{vv'},
\phi_{v'v} \in \mathbb{R}^{K}$ --- one associated with each
endpoint --- and let $\phi$ denote the full collection. Then
\begin{equation}
    \Delta^{\mathrm{LP}}(x, y, w) \;=\;
    \min_{\phi}\, U(x, y, w, \phi) \;-\; \langle w, \psi(x, y) \rangle,
    \label{eq:lp-m3n-loss}
\end{equation}
with
\begin{equation}
    \begin{split}
    U(x, y, w, \phi) \;=\;
    &\sum_{v \in V}\, \max_{a \in \mathcal{Y}}\!
    \Biggl[\,
        f_{v}(x, a) + \tfrac{1}{N}\,\mathbb{I}[a \neq y_{v}]
        - \!\!\!\sum_{v' \in \mathcal{N}(v)}\!\!\!\phi_{vv'}(a)
    \,\Biggr] \\
    +\,
    &\sum_{\{v, v'\} \in E}\, \max_{a, b \in \mathcal{Y}}\!
    \Bigl[\,
        f_{vv'}(a, b) + \phi_{vv'}(a) + \phi_{v'v}(b)
    \,\Bigr].
    \end{split}
    \label{eq:lp-m3n-U}
\end{equation}
For a fixed $\phi$, evaluating $U$ costs
$O(|V| K + |E| K^{2})$: each node and edge term is an
independent local arg-max. No LP solver is invoked.

The inner minimum over $\phi$ in Eq.~\eqref{eq:lp-m3n-loss} is not
solved exactly. Instead, the minimization is folded into the
overall training objective: each training example $i$ carries its
own auxiliary variables $\phi^{i}$, and the per-example duals are
collected into $\phi = (\phi^1, \ldots, \phi^m)$ and optimized
jointly with the model parameters $w$:
\begin{equation}
    (w^{*}, \phi^{*}) \;=\;
    \argmin_{w,\, \phi}\!\left[
    \frac{\lambda}{2}\, \|w\|^{2}
    \;+\; \frac{1}{m}\, \sum_{i=1}^{m}
    \Bigl(
        U(x^i, y^i, w, \phi^{i})
        - \langle w, \psi(x^i, y^i) \rangle
    \Bigr)
    \right].
    \label{eq:lp-m3n-joint}
\end{equation}
Formulation~\eqref{eq:lp-m3n-joint} is the one implemented in this
thesis.

\subsection{Training procedure for LP-M3N}
\label{subsec:training-procedure}

Objective~\eqref{eq:lp-m3n-joint} is a sum over training examples
and contains the neural backbone parameters $\theta$, the pairwise
matrix $W$, and the per-example duals $\phi^{i}$. It is convex in
$(W, \phi)$ for fixed unaries but non-convex in $\theta$, and it is
minimized by mini-batch stochastic optimization: at each step a
small batch $\mathcal{B} \subset \{1, \ldots, m\}$ replaces the
full sum and the parameters $(\theta, W, \phi)$ are updated using
the corresponding subgradient. Plain stochastic gradient descent
(SGD) with a fixed learning rate is brittle on neural networks, so
the Adam optimizer is used throughout. Adam maintains per-parameter
exponential moving averages of the gradient and of its squared
magnitude, yielding adaptive step sizes that empirically give
stable convergence on neural networks of the kind used in this
thesis.

At each training step, a forward pass through the neural backbone
produces the unary potentials $f_{v}(x, \cdot)$; the per-example
terms inside the sum of Eq.~\eqref{eq:lp-m3n-joint} are evaluated
in closed form using the current $\phi^{i}$, and subgradients are
obtained via automatic differentiation of the $\max$ operations
(the subgradient of a max is the indicator at the argmax). The
auxiliary variables $\phi^{i}$ are initialized to zero and stored
persistently across epochs --- one set per training example --- and
are updated jointly with $(\theta, W)$ by the same optimizer. No
inference oracle is called during training; MAP inference
(Section~\ref{sec:map-inference}) is invoked only at evaluation
time.
